\documentclass[journal ]{new-aiaa}
%\documentclass[conf]{new-aiaa} for conference papers
\usepackage[utf8]{inputenc}
\usepackage{textcomp}

\usepackage{graphicx}
\usepackage{amsmath}
\usepackage[version=4]{mhchem}
\usepackage{siunitx}
\usepackage{longtable,tabularx}

\usepackage[most]{tcolorbox}

% Define "problem" environment
\newtcolorbox{problem}[1][]{%
  enhanced,
  colback=white,      % background color
  colframe=black,     % border color
  boxrule=0.8pt,      % border thickness
  arc=2pt,            % corner roundness
  left=6pt, right=6pt, top=6pt, bottom=6pt,
  title=\textbf{Problem: #1},
}



% zaciatok mojho
\let\openbox\relax  % newtxmath (nacitany triedou) uz \openbox definuje; bez tohto amsthm skonci chybou
\usepackage{amsthm}

% Numbering within sections
\numberwithin{equation}{section}

\theoremstyle{definition}
\newtheorem{definition}{Definition}[section]

\theoremstyle{plain}
\newtheorem{theorem}[definition]{Theorem}
\newtheorem{lemma}[definition]{Lemma}
\newtheorem{claim}[definition]{Claim}
\newtheorem{corollary}[definition]{Corollary}
\newtheorem{proposition}[definition]{Proposition}

\theoremstyle{remark}
\newtheorem{remark}[definition]{Remark}
\newtheorem{example}[definition]{Example}


\usepackage{algpseudocode}  % part of algorithmicx package
\usepackage{algorithm}
\usepackage{fancyvrb}       % improved verbatim environment
\usepackage{listings}       % pretty-printer of source code
\usepackage{subcaption}
\let\Bbbk\relax     % to iste pre amssymb: \Bbbk uz definuje newtxmath
\usepackage{amssymb}
\usepackage{etoolbox}
\AtBeginEnvironment{algorithmic}{\setstretch{1}}

%koniec mojho

\setlength\LTleft{0pt} 

\title{Min Cut-Path Problem}

\author{Norbert Micheľ\\
Institute of Computer Science, Faculty of Science, Pavol Jozef Šafárik University in Košice\\
Košice, Slovakia\\
\texttt{5344553@upjs.sk}}


\begin{document}

\maketitle

\begin{abstract}

In modern distributed systems, data exchange between servers must often balance two opposing objectives: 
\emph{enabling communication} within a trusted domain and \emph{preventing communication} by potential adversaries. 
When information flows through a network, it is desirable to establish a reliable path between authorized servers 
while simultaneously enforcing structural constraints that deny any alternative communication routes accessible to an opponent. 
This dual requirement leads naturally to the study of network configurations that combine 
\emph{communication} for legitimate participants with \emph{control} mechanisms that isolate or restrict hostile ones. 
In this context, one seeks structures that ensure communication only within a specified subnetwork while suppressing all competing connections. 
In this paper, we formalize this concept through the notion of a \emph{cut-path}, 
a structure that simultaneously contains a connecting path and a separating cut, 
and we study the computational and structural properties of the corresponding \textsc{Min Cut-Path} problem.

\end{abstract}





\section{Introduction}

In distributed communication systems, a central challenge lies in maintaining a secure but controlled flow of information between servers. 
Although authorized entities require uninterrupted communication channels, adversarial agents must be denied access to any possible route that would allow information transfer. 
This tension between enabling and restricting communication naturally motivates the study of graph structures that can simultaneously express both communication and isolation.
The interplay between these two objectives connectivity and control serves as the conceptual foundation for the present work.

Graph theory provides an elegant formal framework for representing such systems. 
Vertices correspond to servers or agents, and edges represent direct communication links between them. 
Within this model, ensuring communication between trusted servers while cutting off an opponent translates to identifying a subset of edges that forms both a connecting path and a separating cut. 
This hybrid structure, which captures both connection and disconnection properties at the same time, will be formally introduced as a \emph{cut-path}. 
The existence of such structures allows the system to maintain communication within a trusted subnetwork while preventing any external infiltration or alternative routes accessible to an adversary.

Several classical problems in network theory have attempted to combine notions of connectivity and separation from different perspectives. 
For example, the study of \emph{cut-pairs}, \emph{minimum cutsets}, and \emph{network reliability paths}~\cite{GomoryHu1961,NagamochiKameda1996,Mehlhorn2016Certifying} explores trade-offs between maintaining essential communication routes and minimizing potential vulnerabilities. 
Similarly, the interaction between \emph{shortest paths} and \emph{minimum cuts} has been investigated in the context of \emph{maximum flow-minimum cut dualities}~\cite{Cormen2022}. 
However, these formulations treat connectivity and separation as distinct or opposing properties, rather than unifying them within a single combinatorial object. 
In contrast, the present work introduces a unified perspective through the concept of a \emph{cut-path}, which inherently embodies both connectivity and disconnection within the same set of edges.


The central problem studied in this paper, the \textsc{Min Cut-Path} problem, seeks the smallest possible set of edges that simultaneously forms a path and a cut between two distinguished vertices. 
This dual optimization objective reflects the balance between maintaining minimal communication infrastructure and enforcing maximal control over connectivity. 
Despite its natural formulation, the problem has not been previously studied in the literature, and its computational properties remain unexplored. 
We address this gap by providing the first formal analysis of its complexity and structural characteristics.

Our main contributions are as follows. 
Firstly, we prove that the decision version of the \textsc{Min Cut-Path} problem is NP-complete through a two-step reduction from 3-SAT through an intermediate problem called \textsc{Separating Shortest Path}. 
Secondly, we identify important graph classes in which the problem becomes tractable. 
In particular, we show that for graphs with diameter two or with minimum cut-value at most two, the value of the minimum cut-path can be computed in polynomial time and expressed by a simple closed formula. 
Finally, we connect these results with random graph theory, demonstrating that dense Erdős–Rényi graphs almost surely fall into the tractable regime, leading to efficient average-case approximations.

The remainder of this paper is structured as follows. Section~\ref{sec:Fundamentals} introduces the formal definitions, notation, and fundamental concepts used throughout the paper, including the definition of cut-paths and their properties. The NP-completeness of the \textsc{Min Cut-Path} problem is established in Section~\ref{sec:NP-completeness } through a two-step reduction from \textsc{3-SAT} via the intermediate \textsc{Separating Shortest Path} problem. In Sections~\ref{sec:diameter_two} and~\ref{sec:cut_two}, a detailed analysis is presented for the classes of graphs in which the problem becomes polynomially solvable, specifically for graphs of diameter two and graphs with minimum cut-value at most two. The structural properties of Erdős–Rényi random graphs are further investigated in Section~\ref{sec:Diameter}, with emphasis on their diameter, connectivity, and implications for the \textsc{Min Cut-Path} problem. Finally, Section~\ref{sec:conclusion} concludes the paper and provides an average-case approximation result for the \textsc{Min Cut-Path} problem based on the observed structural characteristics.






\section{Fundamentals}
\label{sec:Fundamentals}

Throughout the paper, we use standard notation from graph theory and computational complexity, following the conventions of~\cite{Diestel2023,Cormen2022}. In general, we adopt the convention that all graphs under consideration are undirected and finite, unless otherwise stated. For a graph \(G = (V, E)\), we let \(n = |V|\) denote the number of vertices and \(m = |E|\) denote the number of edges, unless otherwise specified. 

We will often deal with two distinct vertices \(u, v \in V\). The value $d(u, v)$ denotes the distance between \(u\) and \(v\), i.e., the length of shortest \(u\)–\(v\) path in \(G\). Similarly, the value $c(u, v)$ denotes the minimum size of a cut that separates \(u\) and \(v\). All problems studied in this paper are defined with respect to these two distinguished vertices
\(u\) and \(v\). 

For any path \(P \subseteq E\), we denote by \(G \setminus P\) the graph obtained from \(G\) by removing all edges of the path \(P\). Furthermore, for a vertex \(v \in V\) and a vertex subset \(S \subseteq V\), we write \(\deg(v, S)\) for the number of neighbors of \(v\) that belong to \(S\), that is \(\deg(v, S) = |\{\,u \in S : \{u, v\} \in E\,\}|.\)


The notion of a \emph{cut-path} captures a hybrid structure in a graph that combines two fundamental and seemingly opposing elements: a \textit{path} and a \textit{cut}. Intuitively, a \textit{cut-path} between two vertices \(u\) and \(v\) is a set of edges that simultaneously contains a path connecting \(u\) to \(v\), and a cut that, if removed, would disconnect \(u\) from \(v\). In this sense, a cut-path ensures both \textit{communication} and \textit{control}. More Formally:

\begin{definition}[Cut-Path]
Let \(G = (V, E)\) be a graph, and let \(u, v \in V\) be two distinct vertices. A cut-path between \(u\) and \(v\) is a set of edges \(S \subseteq E\) such that, there exist subsets \(P, C \subseteq S\) where:
    \begin{enumerate}
        \item \(C\) is a cut separating \(u\) and \(v\) (i.e., removing \(C\) disconnects \(u\) and \(v\)),
        \item \(P\) is a path connecting \(u\) and \(v\) (i.e., \(P\) forms a sequence of edges between \(u\) and \(v\)).
    \end{enumerate}
\end{definition}

\begin{definition}[Cut-Paths]
Let \(G = (V, E)\) be a graph, and let \(u, v \in V\) be two distinct vertices. The set of all cut-paths between \(u\) and \(v\), denoted by \(\text{CP}(u, v)\), is defined as:
\[
\text{CP}(u, v) = \{ S \mid S \text{ is a cut-path between } u \text{ and } v \}.
\]
\end{definition}

Finally, we use \(\text{cp}(u, v)\) to denote the value of a minimum cut-path between \(u\) and \(v\), 
that is, the minimum number of edges among all sets that are cut-paths.

\begin{definition}[Value of Minimum Cut-Path]
Let \(G = (V, E)\) be a graph, and let \(u, v \in V\) be two distinct vertices. The minimum size of any cut-path between \(u\) and \(v\), denoted by \(\text{cp}(u, v)\), is defined as:
\[
\text{cp}(u, v) = \min\{ |S| \mid S \in \text{CP}(u, v) \}.
\]
\end{definition}

In this paper, we will refer to a \textsc{Min Cut-Path} between vertices \(u\) and \(v\) as any cut-path \(S \in \text{cut-path}(u, v)\) whose cardinality equals the minimum cut-path value \(\text{cp}(u, v)\), i.e., \(|S| = \text{cp}(u, v)\). Formally:

\begin{problem}[\textsc{Min Cut-Path} (Optimization Version)]
Let $G=(V,E)$ be an undirected graph and $u,v \in V$ distinct vertices.  
A set of edges $S \subseteq E$ is called a \emph{cut-path} if it can be partitioned into two subsets $S = P \cup C$ such that $P$ forms a $u$--$v$ path in $G$ and $C$ forms a $u$--$v$ cut in $G$.  

The optimization version of the \textsc{Min Cut-Path} problem is defined as follows:

\begin{center}
\begin{tabular}{p{0.18\linewidth}p{0.75\linewidth}}
\textbf{Input:} & A graph $G=(V,E)$ and vertices $u,v \in V$. \\
\textbf{Output:} & A minimum cut-path $S ^* \subseteq E$ between $u$ and $v$, together with its cardinality
\[
cp(u,v) = |S^*|.
\]
\end{tabular}
\end{center}
\end{problem}


Before formalizing the notion of an average $(1+\epsilon)$-approximation scheme, we briefly introduce the motivation behind this concept. 
In the context of \textsc{Min Cut-Path}, exact computation may be infeasible for large random. 
However, many such graphs exhibit structural regularities that make near-optimal solutions achievable in polynomial time on almost all inputs. 
The following definition captures this idea by requiring that the algorithm performs within a $(1+\epsilon)$ factor of the optimum for almost every instance while still maintaining the polynomial runtime.

To make the definition precise, let $I(n)$ denote the set of all valid input instances of size $n$, and let $\mathrm{OPT}(x)$ represent the optimal value of the objective function for a given instance $x \in I(n)$. 
For minimization problems, $\mathrm{OPT}(x)$ corresponds to the minimum achievable value, while for maximization problems it denotes the maximum achievable value. Formally:

\begin{definition}
[Average $(1+\epsilon)$-Approximation Scheme]
Let $\epsilon > 0$. An algorithm $A$ is called an \emph{average $(1+\epsilon)$-approximation scheme} if it meets the following criteria:
\begin{enumerate}
    \item $A$ is a polynomial-time algorithm, i.e., its running time is $O(P(n))$ for some polynomial $P(n)$.
    \item For each input size $n$, there exists a subset of inputs $I_{A}^{\mathrm{opt}}(n) \subseteq I(n)$ such that for every input $x \in I_{A}^{\mathrm{opt}}(n)$ the output $S = A(x)$ satisfies
    \[
    \frac{S}{\mathrm{OPT}(x)} < 1+\epsilon \quad \text{(for minimization problems)}
    \]
    \[
    \frac{\mathrm{OPT}(x)}{S} < 1+\epsilon \quad \text{(for maximization problems)},
    \]
    where $\mathrm{OPT}(x)$ denotes the optimal value for the input $x$.
    \item In addition, the algorithm is an average-case in the sense that:
    \[
    \lim_{n \to \infty} \frac{\lvert I_{A}^{\mathrm{opt}}(n) \rvert}{\lvert I(n) \rvert} = 1.
    \]
\end{enumerate}
\end{definition}





\section{NP-completeness of the Min Cut-Path Problem}
\label{sec:NP-completeness }
In this section, we show that the \textsc{Min Cut-Path} problem is NP-complete. 
Our proof proceeds via a sequence of polynomial-time reductions. 
We begin with the well-known \textsc{3-SAT} problem, which is formally defined as follows.





\begin{problem}[\textsc{3-SAT}]
Let $U = \{x_1, x_2, \dots, x_n\}$ be a finite set of Boolean variables, and let 
$\mathcal{C} = \{C_1, C_2, \dots, C_m\}$ be a collection of clauses, where each clause 
$C_j$ is a disjunction of exactly three literals, and each literal is either a variable $x_i$ or its negation $\neg x_i$.  

The \textsc{3-SAT} decision problem is defined as follows:

\begin{center}
\begin{tabular}{p{0.12\linewidth}p{0.8\linewidth}}
\textbf{Input:} & A finite set of Boolean variables $U$ and a collection $\mathcal{C}$ of clauses, each containing exactly three literals.\\
\textbf{Question:} & Does there exist a truth assignment $\tau: U \rightarrow \{\mathsf{true}, \mathsf{false}\}$ such that every clause $C_j \in \mathcal{C}$ evaluates to \textsf{true} under $\tau$?
\end{tabular}
\end{center}
\end{problem}


It is a well-established and widely recognized result in computational complexity theory that \textsc{3-SAT} is NP-complete~\cite{cook1971complexity}, a theorem that serves as a 
foundation for a vast number of subsequent NP-completeness proofs.

Our approach is based on a two-step reduction framework. Instead of reducing \textsc{3-SAT} to \textsc{Min Cut-Path} directly, we introduce an intermediate problem, the \textsc{Separating Shortest Path} problem. First, we show its NP-completeness 
by providing a polynomial-time reduction from \textsc{3-SAT} 
(see Section~\ref{Separating_Shortest_Path}). 
In the second step, we reduce \textsc{Separating Shortest Path} to \textsc{Min Cut-Path} 
(see Section~\ref{SSP_to_MCP}). 
This sequence of reductions (from \textsc{3-SAT} to \textsc{Separating Shortest Path}, and from \textsc{Separating Shortest Path} to \textsc{Min Cut-Path}) establishes the NP-completeness of \textsc{Min Cut-Path}.


\subsection{The Separating Shortest Path Problem}\label{Separating_Shortest_Path}
In contrast to the \textsc{Min Cut-Path} problem, where the objective is to find a set of edges that simultaneously contains both a $u$--$v$ path and a $u$--$v$ cut, the \textsc{Separating Shortest Path} problem focuses solely on identifying a shortest path with a specific separation property.  
Namely, we ask whether, among all shortest $u$--$v$ paths in the graph, there exists one whose removal disconnects $u$ and $v$.  
In other words, the path itself is not part of an optimization over combined structures, but rather a candidate from the set of shortest paths satisfying a separation criterion.

\begin{problem}[\textsc{Separating Shortest Path}]
Let $G = (V,E)$ be an undirected graph and let $u,v \in V$ be two distinct vertices.  
A shortest $u$--$v$ path $P$ is called a \textsc{Separating Shortest Path} if the graph $G \setminus P$ contains no $u$--$v$ path.  
The \textsc{Separating Shortest Path} decision problem is defined as follows:

\begin{center}
\begin{tabular}{p{0.12\linewidth}p{0.8\linewidth}}
\textbf{Input:} & A graph $G = (V,E)$ and two distinct vertices $u,v \in V$.\\
\textbf{Question:} & Does there exist a shortest $u$--$v$ path in $G$ that is a separating shortest path?
\end{tabular}
\end{center}
\end{problem}

It is worth noting that, unlike \textsc{Min Cut-Path}, the \textsc{Separating Shortest Path} problem is not an optimization problem; it merely asks for the existence of a shortest path satisfying a certain property.  
This distinction will allow us to construct a more direct reduction from \textsc{3-SAT}.




\subsubsection{Gadget and Basic Operations for the Reduction}

To establish the NP-completeness of \textsc{Separating Shortest Path}, we introduce a collection of \emph{gadgets} and \emph{basic operations} with them.  
These gadgets and operations will be the elementary building blocks of our reduction from \textsc{3-SAT} to \textsc{Separating Shortest Path}.  
They allow us to represent the logical structure within the graph in a modular way.  

The reduction will be illustrated through two parts:
\begin{enumerate}
    \item the \emph{base structure}, also referred to as the \emph{chain};
    \item the \emph{threads}, which will be threading to the base structure.
\end{enumerate}

Figure~\ref{fig:chain_threads} presents an initial schematic overview of the two types of gadgets, illustrating the fundamental layout that underlies the entire construction.

In this setting, computing $cp(u,v)$ becomes significantly more challenging.  
To address this, we use the \emph{general square graph decomposition}, which breaks the graph into overlapping square structures that preserve the essential connectivity and cut properties relevant for min cut-path computations.  
A schematic illustration of this decomposition is shown in Figure~\ref{fig:general_square_graph_pic}.


\begin{figure}[H]
    \centering
    \includegraphics[width=145mm]{img/chain_threads.png}
    \caption{\textbf{Fundamental Layout}}
    \label{fig:chain_threads}
\end{figure}
Having outlined the overall structure, we now turn to a detailed presentation of the individual gadgets and operations that form the core of the reduction.


\paragraph{Chain}  

We begin by introducing the smallest component among the gadgets, referred to as a \emph{chain link}.
This object is a simple cycle that is anchored at two vertices, (p) and (q), and is composed of two distinct paths connecting these vertices, each of equal length and overlapping only at their end-vertices.


\begin{definition}[Chain Link]
Let $G=(V,E)$ be a graph and let $p,q \in V$ be two distinct vertices.  
A \emph{$(p,q)$-chain link} is a subgraph $L \subseteq G$ such that:
\begin{enumerate}
    \item $L$ is a simple cycle in $G$;
    \item the vertices $p$ and $q$ lie on $L$ and partition it into two internally disjoint $p$--$q$ paths;
    \item these two $p$--$q$ paths have equal length.
\end{enumerate}
We refer to one of the two $p$--$q$ paths as the \emph{positive path} of the chain link and the other as the \emph{negative path}. Without loss of generality, when using a chain link inside a chain, its boundary vertices are exactly $p$ and~$q$, corresponding to the points where the adjacent single edges of the chain are attached.
\end{definition}


\begin{figure}[H]
    \centering
    \includegraphics[width=50mm]{img/chain_link.png}
    \caption{\textbf{Chain Link}}
    \label{fig:chain_link}
\end{figure}



The chain forms the backbone of the construction. Its defining feature is the strict alternation between single edges and \emph{chain links}.  
The structure always begins with a single edge, then continues with a chain link, then again a single edge, and so on, and finally ends with a single edge. Formally: 

\begin{definition}[Chain]
Let $G=(V,E)$ be a graph, and let $u,v\in V$.  
We say that $G$ is a \emph{$u$--$v$ chain} if the edge set $E$ can be partitioned into egde-disjoint subgraphs
\[
E = H_1 \cup L_1 \cup H_2 \cup L_2 \cup \cdots \cup H_m,
\]
where:
\begin{enumerate}
    \item each $H_i$ consists of a single edge with end-vertices $p_i,q_i$;
    \item each $L_i$ is a \emph{$q_i,p_{i+1}$-chain link}, i.e., a connected subgraph of $G$ with two distinguished end-vertices $q_i,p_{i+1}$;
    \item the first end-vertex $p_1$ coincides with $u$ and the last end-vertex of $H_m$ coincides with $v$.
\end{enumerate}


\begin{figure}[H]
    \centering
    \includegraphics[width=145mm]{img/chain.png}
    \caption{\textbf{Chain}}
    \label{fig:chain}
\end{figure}


\end{definition}



In the reduction from \textsc{3-SAT} to the \textsc{Separating Shortest Path}, chain links have specific roles that reflect the structure of clauses and variables.  
Let $C_1, C_2, \dots, C_m$ denote the clauses of the input formula and let $x_1, x_2, \dots, x_n$ denote its variables.  
Each chain link is assigned to exactly one of the following three types:

\begin{enumerate}
    \item \textbf{Initialization chain link.}  
    For each variable $x_i$ we introduce an initialization chain link, denoted $I_i$.  
    This link represents the beginning of the traversal corresponding to the variable $x_i$. The total number of initialization chain links is $n$.

    \item \textbf{Terminal chain link.}  
    For each variable $x_i$ we also introduce a terminal chain link, denoted $T_i$.  
    This link marks the end of the traversal associated with the variable $x_i$. The total number of terminal chain links is $n$.

    \item \textbf{Literal chain link.}  
    For each literal $\ell_j$ of clause $C_k$ that corresponds to variable $x_i$, we introduce a literal chain link, denoted $L_{j,k,i}$ or $L_{j,k}$.  
    This link represents the occurrence of the literal $\ell_j$ in the clause $C_k$ with respect to the variable $x_i$.
    The total number of literal chain links is $3m$.
    
\end{enumerate}


\begin{figure}[H]
    \centering
    \includegraphics[width=145mm]{img/Chain_link_example.png}
    \label{fig:Chain_link_example}
\end{figure}

In total, the construction contains exactly $3m + 2n$ chain links, accounting for all initialization, terminal, and literal links.

\paragraph{Threads and Threading.}  
Having defined the chain and its gadgets, we now introduce the notion of \emph{threads}, which represent paths that traverse through the chain links.  
The operation of inserting such paths into the chain will be called \emph{threading}. More precisely:

\begin{definition}[Thread]
Let $G=(V,E)$ be a graph containing a $u$--$v$ chain with chain links $L_1,L_2,\dots,L_m$.  
A \emph{thread} is a $u$--$v$ path $P \subseteq G$ such that:
\begin{enumerate}
    \item $P$ starts at $u$ and ends at $v$;
    \item for each chain link $L_i$, the path $P$ shares with $L_i$ at most one edge;
    \item any two distinct threads $P_1, P_2$ are edge-disjoint and intersect only at the vertices $u$ and $v$.
\end{enumerate}
\end{definition}

Thus, threads can be seen as edge-disjoint $u$--$v$ paths that “pass through” the sequence of chain links, each touching a link in at most one edge.  
The operation of \emph{threading} refers to the process of embedding such threads into the chain.

Threads interact with chain links through the operation of \emph{threading}.  
Intuitively, threading takes an unused edge of a chain link (that has not yet been traversed by any existing thread), subdivides it into three consecutive edges, and then uses the middle edge as the unique crossing edge of the thread through that chain link.  Formally:

\begin{definition}[Threading]
Let $L$ be a chain link in a $u$--$v$ chain, and let $P$ be a thread not yet passing through $L$.  
The \emph{threading operation} on $(P,L)$ consists of the  this elementary operations:
\begin{enumerate}
    \item Select an edge $e=\{x,y\}$ of $L$ that is not used by any existing thread.
    \item Subdivide $e$ into a path $x$--$z_1$--$z_2$--$y$ of three edges, introducing two new vertices $z_1,z_2$.
    \item Extend $P$ so that it traverses the middle edge $\{z_1,z_2\}$ as its unique crossing of the chain link $L$.
\end{enumerate}

\begin{figure}[H]
    \centering
    \includegraphics[width=145mm]{img/Threading.png}
    \caption{\textbf{Threading}}
    \label{fig:Threading}
\end{figure}
If the selected edge $e$ belongs to the positive path of $L$, we say that $P$ is \emph{positively threaded through $L$}; if $e$ belongs to the negative path of $L$, we say that $P$ is \emph{negatively threaded through $L$}.  We say in general that the thread $P$ is \emph{threaded through $L$ at edge $e$}.
\end{definition}


To capture the logical structure of the \textsc{3-SAT} formula, each thread will be assigned a specific interpretational role.  
We distinguish three types of threads:

\begin{enumerate}
    \item \textbf{2-synchronization thread.}  
    This thread connects an initialization chain link $I_i$ with the corresponding terminal chain link $T_i$.  
    Its purpose is to enforce synchronization, ensuring that every separating shortest path must consistently traverse either the positive paths or the negative paths in both $I_i$ and $T_i$.

    \item \textbf{3-synchronization thread.}  
    This thread connects the initialization chain link $I_i$, a literal chain link $L_{j,k,i}$, and the terminal chain link $T_i$ corresponding to the same variable $x_i$.  
    Its role is to synchronize the chosen traversal in the literal chain link with that of the initialization and terminal links for the variable $x_i$.

    \item \textbf{Clause thread.}  
    For each clause $C_k = (\ell_{1} \vee \ell_{2} \vee \ell_{3})$, a clause thread is introduced that traverses the triple of literal chain links corresponding to $\ell_{1}, \ell_{2}, \ell_{3}$.  
    This construction simulates the satisfaction requirement of the clause: in order for a separating shortest path to exist, at least one of the three literals must be realized positively.
\end{enumerate}

\subsubsection{NP-completeness of \textsc{Separating Shortest Path}}
We now prove that \textsc{Separating Shortest Path} is NP-complete by a polynomial-time reduction from \textsc{3-SAT}.  
The reduction is constructive and relies on the chain-exploiting and threading gadgets introduced earlier. Before presenting the full construction, we first introduce a few auxiliary procedures and notational conventions that will allow us to describe the reduction precisely and concisely.

\paragraph{Auxiliary procedures.}
The reduction uses the following auxiliary procedures:
\begin{itemize}
    \item $\textsc{BuildChain}(u,v,m)$: constructs a $u$--$v$ chain with $m$ chain links, where each chain link is realized as a $4$-cycle.
    \item $\textsc{Thread}(u,v,[(L_1,\sigma_1),\dots,(L_k,\sigma_k)])$: creates a thread starting at $u$ and ending at $v$, passing through the chain links $L_1,\dots,L_k$, where $\sigma_i \in \{+,-\}$ specifies whether the threading is positive or negative.
\end{itemize}

\paragraph{Auxiliary notation.}
We also introduce the following definitions and notational conventions:
\begin{itemize}
    \item Literal indices: for a literal $\ell$ we use $\ell[0]$ to denote the clause index in which it occurs and $\ell[1]$ to denote its position inside the clause.
    \item Sign function: for each literal $\ell$ we define
\[
s(\ell) = 
\begin{cases}
+ & \text{if $\ell = x_i$ for some variable $x_i$,}\\
- & \text{if $\ell = \neg x_i$ for some variable $x_i$.}
\end{cases}
\]
\end{itemize}

\begin{theorem}
\label{theorem:separating_shortest_path_np_complete}
The \textsc{Separating Shortest Path} problem is NP-complete.
\end{theorem}

\begin{proof}
 
\paragraph{Reduction from 3-SAT}
Given an instance of \textsc{3-SAT}, we now describe how to construct a graph $G$ with two distinguished vertices $u,v$ such that solving the \textsc{Separating Shortest Path} problem on $(G,u,v)$ corresponds to deciding the satisfiability of the original formula.  
The construction proceeds in polynomial time and is summarized in Algorithm~\ref{alg:reduction}.


\begin{algorithm}[H]
\caption{Reduction from \textsc{3-SAT} to \textsc{Separating Shortest Path}}
\label{alg:reduction}
\begin{algorithmic}[1]
\Require Variables $X=\{x_1,\dots,x_n\}$, clauses $C=\{C_1,\dots,C_m\}$
\Ensure Graph $G$, vertices $u,v$
\State $n \gets |X|$, $m \gets |C|$
\State $G \gets \textsc{BuildChain}(u,v,2n+3m)$
\State $\mathcal{M} \gets$ map variable $\to$ list of literals
\For{$i \gets 1$ to $n$}
    \State $Lits \gets \mathcal{M}[x_i]$
    \State reserve $2+|Lits|$ chain links
    \State label first as $I_i$, last as $T_i$
    \For{each $\ell \in Lits$}
        \State assign next free link as $L_{\ell[0],\ell[1],i}$
    \EndFor
\EndFor
\For{$i \gets 1$ to $n$}
    \State $\textsc{Thread}(u,v,[(I_i,+),(T_i,-)])$
    \State $\textsc{Thread}(u,v,[(I_i,-),(T_i,+)]$
    \State $Lits \gets \mathcal{M}[x_i]$
    \For{each $\ell \in Lits$}
        \State $\textsc{Thread}(u,v,[(I_i,+),(L_{\ell[0],\ell[1],i},-),(T_i,+)])$
        \State $\textsc{Thread}(u,v,[(I_i,-),(L_{\ell[0],\ell[1],i},+),(T_i,-)])$
    \EndFor
\EndFor
\For{$i \gets 1$ to $m$}
    \State $\ell_1,\ell_2,\ell_3 \gets C_i$,
    \State $\textsc{Thread}(u,v,[(L_{\ell_1[0],\ell_1[1]},s(\ell_1),(L_{\ell_2[0],\ell_2[1]},s(\ell_2),(L_{\ell_3[0],\ell_3[1]},s(\ell_3))$
\EndFor
\State \Return $G$
\end{algorithmic}
\end{algorithm}
\paragraph{Correctness Proof}
We now argue that the reduction is correct by showing that the constructed instance of \textsc{Separating Shortest Path} is equivalent to the original \textsc{3-SAT} instance.  
The proof relies on the interpretation of clause threads and synchronization threads, which enforce the logical structure of the formula within the graph.

\medskip
\noindent
\textbf{Equivalence of clause threads.}  
Clause threads ensure that any separating shortest path must intersect at least one literal edge per clause, thereby faithfully simulating the requirement that each clause in the \textsc{3-SAT} formula is satisfied.

\medskip
\noindent

\textbf{Synchronization of variable gadgets.}  
\begin{itemize}
    \item \emph{2-synchronization threads:} enforce consistency between the initialization link $I_i$ and the terminal link $T_i$ of variable $x_i$.  
    \item \emph{3-synchronization threads:} enforce consistency between literal links $L_{j,k,i}$ and the already synchronized pair $(I_i,T_i)$.  
\end{itemize}
Both types of synchronization threads guarantee that the separation respects a consistent truth assignment.

\medskip
\noindent
\medskip
\noindent
\textbf{Two-way correspondence.}  
We complete the proof by establishing the equivalence:
\begin{enumerate}[label=(\alph*)]
    \item If there exists a separating shortest path in $G$, then one can extract a satisfying assignment for the original \textsc{3-SAT} instance.  

    Thanks to the synchronization threads, every separating shortest path enforces a consistent choice between the positive and negative sides of each variable gadget.  
    This choice directly determines the truth value of the corresponding variable in the assignment.  
    
    Moreover, each clause thread guarantees that the separating shortest path must traverse at least one literal edge associated with the clause, thereby ensuring that the clause is satisfied.  
    Since the path must be separating, it necessarily intersects every clause thread, which implies that all clauses $C \in \mathcal{C}$ are satisfied simultaneously.  
    Thus, the induced assignment is a satisfying assignment for the original formula.
    
    \item If there exists a satisfying assignment for the \textsc{3-SAT} instance, then one can construct a separating shortest path in $G$.  

Given a satisfying assignment $\tau$, we build a $u$--$v$ path $P$ as follows: for each variable $x_i$, traverse $I_i$, $L_{j,k,i}$  and $T_i$ consistently with the sign prescribed by $\tau(x_i)$ (positive if $\tau(x_i)=\mathsf{true}$, negative if $\tau(x_i)=\mathsf{false}$)

It remains to show that $P$ is separating, i.e., that $G\setminus P$ contains no $u$--$v$ path. Consider any walk starting at $u$ in $G\setminus P$ and analyze all ways it could try to progress:
\begin{enumerate}
  \item \emph{Using a base-chain edge.} Every base-chain edge used by $P$ is removed in $G\setminus P$, hence this option is immediately blocked.
  \item \emph{Using a synchronization thread.} Such a thread first reaches $I_i$. From there, the next move would either (i) attempt to follow the sign taken by $P$ in $(I_i,T_i)$, which is impossible since the unique threaded edge of that sign is already in $P$, or (ii) switch to the opposite sign via a 2- or 3-synchronization thread. By construction of the synchronization threads, any such switch forces an encounter with a threaded edge already used by $P$ (the “consistency lock”), hence this branch is also blocked in $G\setminus P$.
  \item \emph{Using a clause thread.} Entering a clause thread leads to one of its three literal links. In each clause $C_k$ we chose a satisfied literal $\ell_{k,t}$ and threaded $P$ through $L_{k,t,i}$ with sign $s(\ell_{k,t})$. Therefore, every clause thread contains at least one link whose unique crossing edge is in $P$, so any traversal of the clause gadget must meet a removed edge. Attempts to detour via a 3-synchronization thread necessarily encounter the already-threaded opposite-sign edge (by the same consistency lock as above).
\end{enumerate}
Since all possible continuations from $u$ are blocked in $G\setminus P$, there is no $u$--$v$ path after removing $P$; hence $P$ is separating.
\end{enumerate}

\textbf{Polynomial complexity of the reduction.}

It remains to prove show that the reduction runs in polynomial time.  
The initialization of the base chain requires $O(n+m)$ chain links, which is linear in the size of the input.  
All subsequent operations (\textsc{Thread} calls, assignment of labels, and synchronization gadgets) are implemented by constant-size local modifications of the graph and are repeated at most polynomially many times in $n+m$.  
Hence, the overall construction is polynomial in the input size.

\smallskip
\noindent
One technical detail is the calibration of thread lengths: in order to guarantee that all threads yield shortest $u$--$v$ paths of equal length, we must adjust the number of edges in each threaded segment.  
This can be done by padding each gap with a sufficient number of edges, specifically matching the total number of edges in the chain.  
Since the base chain has $O(n+m)$ edges, this adjustment requires at most $O((n+m)^2)$ edges in total.  
As the number of gaps is polynomial, the construction remains polynomially bounded.

\smallskip
\noindent
\textbf{Polynomial verification.}  
Given a candidate separating shortest path $P$ in $G$, one can verify its validity in polynomial time:  
first check that $P$ is indeed a shortest $u$--$v$ path, and then remove $P$ from $G$ and run a standard reachability test (e.g., BFS or DFS) to determine whether a $u$--$v$ path remains.  
Both steps are computable in polynomial time.
   
\end{proof}

\subsection{From Separating Shortest Path to Min Cut-Path}\label{SSP_to_MCP}
In the previous subsection we proved that \textsc{Separating Shortest Path} is NP-complete.  
We now exploit this result to show that our target problem, \textsc{Min Cut-Path}, is also NP-complete.  
For this purpose, we express it in decision form.

\begin{problem}[\textsc{Min Cut-Path}]
Let $G=(V,E)$ be an undirected graph and $u,v \in V$ distinct vertices.  
A set of edges $F \subseteq E$ is called a \emph{cut-path} if $F$ can be partitioned into two subsets $F = P \cup C$ such that $P$ forms a $u$--$v$ path in $G$ and $C$ forms a $u$--$v$ cut in $G$.  
The decision version of the \textsc{Min Cut-Path} problem is defined as follows:

\begin{center}
\begin{tabular}{p{0.12\linewidth}p{0.8\linewidth}}
\textbf{Input:} & A graph $G=(V,E)$, vertices $u,v \in V$, and an integer $k$. \\
\textbf{Question:} & Does there exist a cut-path $F$ with $|F|\leq k$?
\end{tabular}
\end{center}
\end{problem}

We first show that the decision version of the \textsc{Min Cut-Path} problem is NP-complete.  
As an immediate consequence, the corresponding optimization problem is also NP-complete.

\begin{theorem}
\label{theorem:min_cut_path_np_complete}
The decision version of the \textsc{Min Cut-Path} problem is NP-complete. 
\end{theorem}

\begin{proof}
  

\textbf{Reduction Construction.}
Let $G$ be an instance of \textsc{Separating Shortest Path} with distinguished vertices $u,v$.  
We define $G' := G$ and set the threshold $k := d_G(u,v)$, the length of a shortest $u$--$v$ path in $G$.  
We then ask whether $G'$ admits a cut-path of size at most $k$.

\textbf{Correctness of the Reduction.}
Suppose $G$ admits a separating shortest path $P$ of length $d_G(u,v)$.  
Then $P$ itself is a cut-path of size $k$, since (i) it is a $u$--$v$ path and (ii) its removal disconnects $u$ and $v$, forming a $u$--$v$ cut.  
Conversely, if $G'$ has a cut-path $F$ of size $k$, then $F$ must contain a shortest $u$--$v$ path in $G$ and simultaneously separate $u$ from $v$, which is precisely a separating shortest path.  
Thus, the two problems are equivalent under this reduction.

\textbf{Polynomial verification.}  
A candidate cut-path $F$ can be verified in polynomial time:  
check that $F$ contains a $u$--$v$ path, and verify that $G\setminus F$ disconnects $u$ and $v$ (e.g., by a BFS/DFS).  
Both conditions are testable in polynomial time.

\textbf{Conclusion of NP-completeness.} 
We have shown that \textsc{Separating Shortest Path} is NP-complete and that it reduces to \textsc{Min Cut-Path} in polynomial time. Since \textsc{Min Cut-Path} is clearly in NP (by polynomial verification), it follows that \textsc{Min Cut-Path} is NP-complete.
  
\end{proof}

\section{Graphs of Diameter Two or Cut-Value at Most Two}
\label{sec:diameter_two}



Having established the general complexities of the Min Cut-Path problem, we now turn our attention to specific graph classes where we might uncover more tractable solutions. A natural first step in this direction is to consider a constraint on the maximum distance between vertices. Limiting the shortest path length to a maximum of two edges leads directly to the well-known class of graphs with diameter two, which have been extensively studied in algebraic graph theory \cite{GodsilRoyle2001}. We formally define this class as follows:

\begin{definition}
A graph $G = (V, E)$ has diameter $k$ if for every pair of vertices $x, y \in V$, the distance between $x$ and $y$ is at most $k$. Formally:
\[
\forall x, y \in V: \, d(x, y) \leq k,
\]
where $d(x, y)$ denotes the length of a shortest path between $x$ and $y$.

\end{definition}

Our objective in this section is to analyze the structural properties of graphs with diameter two and to utilize these features to derive key theorems. These theorems will then serve as the foundation for developing a straightforward and efficient algorithm for solving the Min Cut-Path problem within this specific class of graphs.  As a first exploration of the structural characteristics of these graphs, we begin by demonstrating a fundamental decomposition property.  We will show that, in any graph, a cut $C$ naturally induces a partition of the set of vertices into subsets $I, J, K, L$, providing a basis for further structural analysis.



The following observation follows directly from the definition of a cut.  
Each vertex either has edges belonging to the cut or not, and it belongs to one of the two sides \(A_1, A_2\).  
Hence, the classification of vertices into subsets \(I, J, K, L\) is immediate and unambiguous:

\begin{lemma}
\label{claim:cut_decomposition}
Let \(G = (V, E)\) be a graph, and let \(C\) be a cut that divides \(G\) into two sets \(A_1\) and \(A_2\) (i.e., \(V = A_1 \cup A_2\) and \(A_1 \cap A_2 = \emptyset\)).  
Then, there exists a unique decomposition of \(A_1\) and \(A_2\) into subsets \(I, J, K, L\) such that:
\begin{itemize}
    \item \(I \cup J = A_1\) and \(I \cap J = \emptyset\),
    \item \(K \cup L = A_2\) and \(K \cap L = \emptyset\),
    \item \(\forall u \in J \cup K, \, \exists v \in V: (u, v) \in C,\)
    \item \(\forall u \in I \cup L, \, \nexists v \in V: (u, v) \in C.\)
\end{itemize}
\end{lemma}


With the introduced vertex decomposition established for any graph and any cut, we now leverage the specific constraint of diameter two to uncover additional structural properties inherent to this class of graphs.  
In particular, for graphs with diameter two, the decomposition induced by a cut \(C\) is no longer arbitrary:  
the limited pairwise distances impose strict constraints on how the vertex subsets \(I, J, K, L\) may coexist.

The key observation is that if both \(I\) and \(L\) were non-empty, then any vertex in \(I\) would be at distance at least three from any vertex in \(L\), since all connections between these two sets must pass through the intermediate subsets \(J\) and \(K\).  
This directly contradicts the assumption that the graph has diameter two.  
Consequently, one of these sets must be empty. Formally:


\begin{lemma}
\label{claim:empty_I_or_J}
Let \(G = (V, E)\) be a graph with diameter 2, and let \(C\) be a cut that divides \(G\) into two sets \(A_1\) and \(A_2\) (i.e., \(V = A_1 \cup A_2\) and \(A_1 \cap A_2 = \emptyset\)). Let \(I, J, K, L\) be the decomposition of \(A_1\) and \(A_2\) as described in Claim~\ref{claim:cut_decomposition}. Then, at least one of the sets \(I\) or \(L\) is empty. Formally:
\[
I = \emptyset \lor L = \emptyset.
\]
\end{lemma}

Without loss of generality, and for the sake of simplifying the analysis, we assume that \(L = \emptyset\).

Now, before focusing again on diameter two graphs, we introduce a lemma outlining a key relationship concerning the intersection of the cut and path within any cut-path, a relationship fundamental to our investigation of the Min Cut-Path problem.

\begin{lemma}
\label{claim:odd_diam_2}
Let \(G = (V, E)\) be a graph, and let \(u, v \in V\) be two distinct vertices. Let \(c(u, v)\) be a cut that separates \(u\) and \(v\), and let \(p(u, v)\) be a path that connects \(u\) and \(v\). Then, the intersection of \(c(u, v)\) and \(p(u, v)\) has an odd number of edges. Formally:
\[
|c(u, v) \cap p(u, v)| \equiv 1 \pmod{2}.
\]
\end{lemma}

\begin{proof}
Let \(A_1\) and \(A_2\) be the two sets defined by the cut \(c(u, v)\), where \(u \in A_1\) and \(v \in A_2\). Consider the path \(p(u, v)\) between \(u\) and \(v\). Let \(S = c(u, v) \cap p(u, v)\) be the set of edges at the intersection of the cut and the path. 

Each edge in \(S\) changes the side of the cut (i.e., switches between \(A_1\) and \(A_2\)) when traversing the path. Without loss of generality, assume that \(u \in A_1\) and \(v \in A_2\). To move from \(A_1\) to \(A_2\), the number of times the path crosses the cut must be odd. Specifically:
\begin{enumerate}
    \item If the path crosses the cut an even number of times, the starting and ending sets would be the same (i.e., \(u\) and \(v\) would both be in \(A_1\) or both in \(A_2\)), which contradicts the assumption that \(c(u, v)\) separates \(u\) and \(v\).
    \item Therefore, the path must cross the cut an odd number of times, i.e., \(|S| \equiv 1 \pmod{2}\).
\end{enumerate}
Thus, the intersection \(c(u, v) \cap p(u, v)\) has an odd number of edges and it completes the proof.
\end{proof}

Building on the structural insights gained from the previous lemmas, particularly the lemmas~\ref{claim:empty_I_or_J} and~\ref{claim:odd_diam_2}, we are now ready to derive a precise formula for the min cut-path value $cp(u,v)$ in graphs with diameter two. Thanks to the specific structure of graphs with diameter two, and using the established relationships between cuts and paths within the cut-path problem, we can express $cp(u,v)$ directly in terms of the minimum cut value $c(u,v)$ and the shortest path distance $d(u,v)$. We are now ready to state the central theorem of this section:

\begin{theorem}
\label{theorem:min_cut_path_diameter_2}
Let \(G = (V, E)\) be a connected graph with diameter 2, and let \(u, v \in V\) be two distinct vertices. Then, the value of the min cut-path between \(u\) and \(v\) satisfies the following:
\[
cp(u, v) = c(u, v) + d(u, v) - 1.
\]
\end{theorem}


\begin{proof}

We distinguish two cases based on the values of \(d(u, v)\) and \(c(u, v)\).

\begin{enumerate}
    \item \textbf{Case} \(d(u, v) = 1\) \textbf{or} \(c(u, v) = 1\): 
    
    It is straightforward to observe that if \(d(u, v) = 1\) or \(c(u, v) = 1\), then the path and the cut necessarily share at least one common edge, yielding
    \[
    cp(u, v) = c(u, v) + d(u, v) - 1.
    \]

    \item \textbf{Case} \(d(u, v) = 2\) \textbf{and} \(c(u, v) \geq 2\):
    
    From a simple observation, we can see that the value of \(cp(u, v)\) must lie between the minimum cut size (or shortest path) and the sum of the cut and path sizes reduced by one, that is,
    \[
    \max(c(u, v), d(u, v)) \leq cp(u, v) \leq c(u, v) + d(u, v) - 1.
    \]
    Consequently,
    \[
    c(u, v) \leq cp(u, v) \leq c(u, v) + 2 - 1 = c(u, v) + 1.
    \]

Assume for contradiction that the equation \(cp(u, v) = c(u, v) + d(u, v) - 1 = c(u, v)+1\) does not hold. Then, the min cut-path problem would have the same size as the min cut, i.e., \(cp(u, v) = c(u, v)\). This implies that the path between \(u\) and \(v\) is a subset of the minimum cut.

By lemma~\ref{claim:odd_diam_2}, the length of the path in the min cut-path problem must be at least 3. By lemma~\ref{claim:empty_I_or_J} and w.l.o.g., the graph can be partitioned into three sets \(I, J, K\), where \(u \in K\) and \(L = \emptyset\). We observe that \(v \in J\). Since there exists a path between \(u\) and \(v\) in the minimum cut-path, and because this path must lie within the minimum cut, it follows that vertex \(v\) cannot belong to the set \(I\) because the path is a subset of the cut. It follows that \( v \in J \). For a better understanding, we provide an illustrative example (see Figure~\ref{fig:2_diam_2}):

\begin{figure}[H]
    \centering
    \includegraphics[width=60mm]{img/2_diam_2.png}
    \caption{\textbf{Illustration of Structure Cut-Path}}
    \label{fig:2_diam_2}
\end{figure}



We now estimate \(cp(u, v)\):
\[
cp(u, v) = c(u,v)\geq \deg(u, I) + \deg(u, J) + 2 + (|K| - 2),
\]
where the estimation is carried out by decomposing the analysis across the vertices:
\begin{itemize}
    \item \(\deg(u, I)\) and \(\deg(u, J)\) denote the number of edges incident from vertex \(u\) to the opposite side of the minimum cut,
    
    \item  from lemma~\ref{claim:odd_diam_2} with \(d(u,v)=2\), any \(u\)-\(v\) path contained in the cut must have length at least three. It follows that there exists at least one vertex other than \(u\) from the set \(K\) that contributes at least two edges to the cut,

    \item the remaining \(|K| - 2\) vertices from the set \(K\) must each contribute at least one edge to the cut.
\end{itemize}


The right-hand side of the inequality is illustrated in red in Figure~\ref{fig:2_diam_2}. 

Simplifying, we obtain:
\[
\deg(u, I) + \deg(u, J) + 2 + (|K| - 2) \geq 
\]
\[
0 +deg(u, J)  + \deg(u, K) + 1 = \deg(u) + 1 \geq c(u, v) + 1.
\]
This contradicts the assumption that \(cp(u, v) = c(u, v)\), proving that the equation \(cp(u, v) = c(u, v) + d(u, v) - 1\) must hold.
\end{enumerate}
\end{proof}



\section{Graphs with Cut-Value at Most Two}
\label{sec:cut_two}

Having thoroughly examined the implications of diameter constraints on the Min Cut-Path problem, we now direct our attention to another significant structural property that offers potential for algorithmic simplification. Specifically, we now focus on the case where the minimum cut between any two vertices is bounded by two. This condition restricts the connectivity structure of the graph in a very particular way.

To clarify the reasoning, we first outline a few basic structural observations. By combining these insights, we immediately obtain the following result (Theorem~\ref{thm:min_cut_path_max_2}):


\begin{enumerate}

    \item \textbf{Step: Cactus structure.}
    Graphs with minimum cut-value at most two necessarily have the structure of a \emph{cactus graph}\cite{Mehlhorn2016Certifying}: they consist of cycles and single edges connected only at articulation points. If two cycles shared more than one independent connection, then separating certain pairs of vertices would require at least three edges, contradicting the assumption \(c(x,y) \leq 2\).
    
    \item \textbf{Step: Formula for \(cp(u,v)\).}  
    Consider a shortest path between two vertices \(u\) and \(v\). In a cactus graph, removing this path still leaves alternative connections within the cycles it traverses. Hence, the cut cannot be fully contained inside the path. It follows immediately that the min cut-path requires exactly one additional edge, yielding the formula:
\[
cp(u,v) = c(u,v) + d(u,v) - 1.
\]
\end{enumerate}



\begin{theorem}
\label{thm:min_cut_path_max_2}
Let \(G = (V, E)\) be a connected graph such that for every pair of distinct vertices \(x, y \in V\), the minimum cut satisfies 
\[
c(x, y) \leq 2.
\]  
Then, for any two vertices \(u,v \in V\), the value of the min cut-path satisfies
\[
cp(u, v) = c(u, v) + d(u, v) - 1.
\]
\end{theorem}









\section{Erdős-Rényi Graphs}
\label{sec:Diameter}

As a first step in our analysis, we examine the structural properties of dense random graphs, specifically those with a high edge probability \( p \). A well-known result in the theory of random graphs states that an Erdős–Rényi random graph \( G(n, p) \) with \( p = \frac{1}{\alpha} \) (where \( \alpha > 0 \) is a constant) almost surely has diameter 2 in the asymptotic limit \cite{Bollobas2001}. 

In the context of the \textsc{Min Cut-Path} problem, this property is particularly insightful, as it allows us to take advantage of the efficient methods developed in Chapter 3 Section \ref{sec:diameter_two} to solve the problem in graphs of diameter 2. One of the key results enabling this is Theorem~\ref{theorem:min_cut_path_diameter_2}, which shows that when the graph has diameter 2, the value of the \textit{min cut-path} between any two vertices \(u\) and \(v\) can be computed in polynomial time via the simple formula:
\[
cp(u, v) = c(u, v) + d(u, v) - 1.
\]

This characterization provides an immediate and efficient way to compute the min cut-path in such graphs, making it a powerful tool for average-case analysis in dense random graphs.



\begin{theorem}
\label{theorem:diameter_2_random_graph}
Let \(G(n, p(n))\) be an Erdős–Rényi random graph with \(n\) vertices and edge probability \(p(n) = \frac{1}{\alpha}\), where \(\alpha > 0\) is a constant. Then, the probability that \(G(n, p(n))\) has diameter 2 converges to 1 as \(n \to \infty\). Formally:
\[
\lim_{n \to \infty} \mathbb{P}\!\left[\text{diam}\bigl(G\bigl(n, \frac{1}{\alpha}\bigr)\bigr) = 2\right] = 1.
\]

\end{theorem}


This chapter investigates the structural properties of the sparse Erdős–Rényi random graphs \( G(n, p) \) with \( p = \frac{\alpha \log n}{n} \) (where \( \alpha > 0 \) is a constant). The primary focus is on fundamental results concerning the diameter and the minimum cut of such graphs. Although these results are well-established in the literatures \cite{Bollobas2001,Frieze2016}, this chapter provides a concise overview, highlighting the statements relevant to the analysis of the \textsc{Min Cut-Path} problem. To maintain clarity and brevity, only the necessary theorems are stated, with proofs omitted except for one result regarding edge connectivity \( c(u,v) \) between arbitrary pairs of vertex.


\begin{theorem} [Diameter~\cite{Bollobas2001}]
\label{theorem:diameter_log}
Let \( G(n, p(n)) \) be an Erdős–Rényi random graph with \( p(n) = \frac{\alpha \log n}{n} \), where \( \alpha > 1 \). Then:
\[
\lim_{n\to\infty} \mathbb{P}\left[ \text{diameter of } G \leq \frac{\log n}{\log \log n} \right] =1.
\]
\end{theorem}


\begin{theorem} [Connectivity~\cite{Frieze2016}]
\label{theorem:k-edge-connectivity}
Let \( G(n, p(n)) \) be an Erdős–Rényi random graph with \( p(n) = \frac{\alpha \log n}{n} \), where \( \alpha > 1 \). Then:
\[
\lim_{n\to\infty} \mathbb{P}\left[ G \text{ is } k\text{-edge-connected} \right] = 1,
\]
where \( k= \delta(G) \) denotes the minimum degree of \( G \).
\end{theorem}



\begin{theorem}[Degree Concentration~\cite{Bollobas2001}]
\label{theorem:degree_concentration}
Let \( G(n, p(n)) \) be an Erdős–Rényi random graph with
\[
p(n) = \frac{\alpha \log n}{n},
\]
where \( \alpha > 1 \) is a fixed constant. Then, for any fixed \( \epsilon > 0 \), it holds that
\[
\lim_{n\to\infty} \mathbb{P}\Bigl( \forall v \in V(G): \ (1 - \epsilon)\alpha \log n \leq \deg(v) \leq (1 + \epsilon)\alpha \log n \Bigr) = 1.
\]
\end{theorem}


The preceding results provide structural constraints on connectivity or the diameter of the graph. However, we must make these findings applicable in the context of the Min Cut-Path problem. The following lemma establishes bounds on the value of \( c(u,v) \):

\begin{lemma}
\label{claim:max_degree}
Let \( G(n, p(n)) \) be an Erdős–Rényi random graph with \( p(n) = \frac{\alpha \log n}{n} \), where \( \alpha > 1 \). Consider two randomly chosen vertices \( u, v \in V(G) \). For every \( \epsilon > 0 \), there exists \( n_0 \) such that for all \( n > n_0 \), the edge connectivity between \( u \) and \( v \) satisfies:
\[
\lim_{n\to\infty} \mathbb{P}\Bigl( (1 - \epsilon)\alpha \log n \leq c(v,v) \leq (1 + \epsilon)\alpha \log n \Bigr) = 1.
\]
\end{lemma}


\begin{proof}
First, we fix the values of \(k\) and \(\epsilon\) as in the proofs of Theorems~\ref{theorem:k-edge-connectivity} and \ref{theorem:degree_concentration}, we know that the graph \( G(n, p) \) is \( k \)-edge-connected with high probability, where \(\ (1 - \epsilon)\alpha \log n \leq k \leq (1 + \epsilon)\alpha \log n  \). This implies that for all pairs of vertices \( u, v \), the edge connectivity satisfies:
\[
c(u, v) \geq \min_{x,y \in V(G)} c(x, y) \geq \min_{z \in V(G)} \deg(z) \geq  (1 - \epsilon)\alpha \log n.
\]


Since the edge connectivity is determined by the minimum number of edges that must be removed to disconnect two vertices, we observe that providing a specific cut may not yield the minimum possible value. However, it gives an upper bound. In particular, for the given pair \( (u,v) \), the degree constraint implies
\[
c(u,v)\leq deg(u) \leq (1 + \epsilon)\alpha \log n \Bigr.
\]

By combining these two inequalities, we obtain the desired result:
\[
\ (1 - \epsilon)\alpha \log n \leq c(u,v) \leq (1 + \epsilon)\alpha \log n \Bigr.
\]
\end{proof}
    


Based on known results on the diameter and connectivity of Erdős–Rényi random graphs, we can estimate the expected deviation from the optimal solution. Specifically, the combination of the minimum cut and shortest path solutions leads to an approximation scheme. Using the diameter bound from Theorem~\ref{theorem:diameter_log} and the connectivity concentration from lemma~\ref{claim:max_degree}, we ensure that the approximation ratio remains within a factor of \( (1+\epsilon) \) in average-case. This guarantees that the proposed approach qualifies as an average \( (1+\epsilon) \)-Approximation Scheme.

\begin{theorem}
Let \( G(n, p) \) be an Erdős–Rényi random graph with \( p \geq \frac{\alpha \log n}{n} \), where \( \alpha > 1 \). Then the algorithm that combines the solutions for the \textsc{Min Cut} and the \textsc{Min Path} problem is an Average \( (1+\epsilon) \)-Approximation Scheme.
\end{theorem}
 
\begin{proof}
If the probability \( p \) increases, the length of the shortest path decreases. From Theorem~\ref{theorem:diameter_log}, we obtain the bound:
\[
d(u,v) \leq \text{diameter of } G\leq \frac{\log n}{\log \log n}.
\]
Using the lemma~\ref{claim:max_degree}, we can bound the edge connectivity:
\[
c(u,v) > (1 - \epsilon')\alpha \log n.
\]
It is easy to see that the minimum cut is no larger than the minimum cut-path:
\[
c(u,v) \leq cp(u,v).
\]

By combining all the inequalities and performing successive substitutions, we obtain:

$$
\frac{S(x)}{\text{OPT}(x)} \leq \frac{c(u,v)+d(u,v)}{c(u,v)}\leq 
\frac{(1 - \epsilon')\alpha \log n + \frac{\log n}{\log \log n}}{(1 - \epsilon')\alpha \log n}=
$$
$$
1+\frac{\frac{\log n}{\log \log n}}{(1 - \epsilon')\alpha \log n}\leq 1 +\epsilon.
$$
Thus, we conclude that the algorithm achieves a Average $(1+\epsilon)$-Approximation Scheme.
\end{proof}


\section{Conclusion}
\label{sec:conclusion}

In this paper, we introduced and studied the \emph{Min Cut-Path} problem, a novel combinatorial structure that simultaneously captures communication and control within a single set of edges. We established the NP-completeness of the decision version of the problem by a two-step reduction from 3-SAT through the intermediate \emph{Separating Shortest Path} problem. Furthermore, we identified several important classes of graphs in which the problem becomes polynomially solvable, including graphs of diameter two and graphs whose minimum cut-value does not exceed two. For these cases, we derived a closed formula for the value of the minimum cut-path, showing that \(cp(u, v) = c(u, v) + d(u, v) - 1.\)
Finally, connecting these findings with the structural properties of random graphs, we showed that dense Erdős–Rényi graphs almost surely fall within this tractable regime, which leads to efficient average-case approximations and provides a bridge between worst-case and probabilistic analysis.

Since the concept of a cut-path is new, there remain numerous promising directions for future research. One particularly appealing direction is the identification of additional \emph{islands of polynomial-time solvability}. Expanding the frontier of tractable cases can help better understand the structural boundaries between efficiently solvable and intractable instances of the \textsc{Min Cut-Path} problem.

A natural starting point in this search is the class of graphs in which every pair of vertices has either a distance two or a minimum cut of size two. This class represents a unifying generalization of the two polynomial cases already identified in this work graphs of diameter two and graphs with minimum cut value at most two. By combining these two structural properties, these two polynomial islands could be merged into a single broader class of graphs where the \textsc{Min Cut-Path} problem could remain tractable. 


Another promising research direction concerns \emph{planar graphs}. The motivation arises from the inherent duality between paths and cuts: in planar graphs, every edge cut in the primal graph corresponds to a cycle in the dual graph and vice versa. This dual relationship suggests that the \textsc{Min Cut-Path} problem might admit more efficient formulations or reductions when restricted to planar graphs. Moreover, several classical problems involving cuts and flows are known to be solvable more efficiently in planar graphs than in general ones, particularly due to the existence of planar duality and specialized algorithms for planar embeddings~\cite{itai1979maximum,henzinger1997faster}. Therefore, planar graphs form a natural candidate for further exploration of polynomial solvability.

Beyond theoretical developments, another potential avenue for progress lies in experimental and empirical analysis. As the present work is purely theoretical, empirical evaluation could provide valuable insight into the average-case performance of algorithms for the \textsc{Min Cut-Path} problem. Studying the behavior of the problem on real-world networks or synthetic random instances could reveal structural patterns that are not captured by worst-case analysis and could guide the design of practical heuristics or approximation methods.

Lastly, a deeper exploration of the structural behavior of graphs with fixed cut-path values could reveal additional regularities. In particular, analyzing graphs with a fixed value of \(cp(u,v)\) can reveal structural principles that further delineate the boundary between polynomially solvable and NP-hard instances, enriching the understanding of this newly introduced combinatorial framework.

\section*{Acknowledgment}

The author would like to express his sincere gratitude to \textbf{Prof.~RNDr.~Martin Loebl,~CSc.}, who proposed the original idea of the \emph{Min Cut-Path} problem and provided essential guidance and valuable discussions that greatly influenced the development and results of this work.


\bibliography{references}

\end{document}
